%!TEX root = TheoreticalPhysicsC.tex


\section{Lecture: The Canonical Ensemble}
Consider a system in contact with its environment (a ``heat bath''), so that energy can be exchanged with it. Suppose the system is modelled with Hilbert space $\mathcal{H}$ and Hamiltonian $\widehat{H}$. The exchange of energy can be modelled mathematically by letting the energy of the system have a fixed mean value $E = \langle \widehat{H}\rangle$.

This situation is the normal case: the electromagnetic field in this lecture hall is in a thermal state, and every system in the room interacts with it. We could not determine the true state even in principle; what we set up is a \emph{hypothesis} with the same expectation values.

We are looking for a state $\rho$ which maximises the entropy under the constraint that the mean value of the energy is held constant. Following the principle of maximum entropy, this problem is described mathematically as follows:
\begin{equation}
	\begin{split}
		\max_{\rho} &\,S(\rho) \\
		\text{subject to} &\, E=\tr(\rho \widehat{H}).
	\end{split}
\end{equation}

Two constraints are still missing here: for $\rho$ to describe a state at all, we must have $\tr(\rho)=1$ and $\rho \ge 0$. In full, the task therefore reads
\begin{equation}
	\max_{\rho} S(\rho) \quad\text{subject to}\quad E=\tr(\rho \widehat{H}),\quad \tr(\rho)=1,\quad \rho\ge 0.
\end{equation}
We are looking for a matrix with $d^2$ degrees of freedom, and $S$ is non-linear in these entries; that a closed-form solution exists at all is a piece of luck.

Such maximisation problems can be solved with the help of Lagrange multipliers: 
\begin{equation}
	\max_{\rho} f(\rho) = -k\tr(\rho\log(\rho)) + \beta \tr(\rho \widehat{H})
\end{equation}
Each constraint has its own multiplier; above, only the one for the energy is written down. Normalisation contributes the term $\gamma\tr(\rho)$, which returns in a moment as $\log(\mathcal{Z})\mathbb{I}$; positivity we check at the end.

\begin{greybox}[Lagrange multipliers]
	Suppose we want the maximum of $f(x,y)$ subject to the constraint $g(x,y)=0$. We introduce a constant $\lambda$ and maximise instead the unconstrained function $F(x,y,\lambda) = f(x,y) + \lambda g(x,y)$; the solution still depends on $\lambda$, which we determine at the end from the constraint. For intuition, think of $\lambda$ as very large, and of its sign as chosen so that a violation is penalised: the maximum can then only lie where the constraint holds.
\end{greybox}

We write the solution as a function of $\beta$: $\rho \equiv \rho(\beta)$. Afterwards we look for the $\beta$ for which the condition $E=\tr(\rho \widehat{H})$ is satisfied. 

To solve the maximisation problem, the variation of $f(\rho)$ must be stationary with respect to all possible variations $\delta \rho$ of $\rho$:
\begin{equation}
	0= f(\rho + \delta\rho) -f(\rho)
\end{equation}
One can show that the function attains its maximal value only if
\begin{equation}
	\log(\rho) + \beta \widehat{H} + \log(\mathcal{Z})\mathbb{I} = 0,
\end{equation}
where $\mathcal{Z}$ is a constant.
The solution of this equation is therefore given by an exponential function:
\begin{equation}
	\rho \propto e^{-\beta \widehat{H}}.
\end{equation}
This step is allowed because $\log(\mathcal{Z})$ was multiplied by the identity: $\mathbb{I}$ commutes with everything, and only for that reason does $e^{-\beta \widehat{H}-\log(\mathcal{Z})\mathbb{I}}$ factorise into a product.
The constant of proportionality is precisely $\mathcal{Z}$, and so we obtain
\begin{equation}
	\rho = \frac{e^{-\beta \widehat{H}}}{\mathcal{Z}},
\end{equation}
where the condition $\tr(\rho) = 1$ means that the \emph{partition function} $\mathcal{Z}$ is given by
\begin{equation}
	\mathcal{Z} = \tr\left(e^{-\beta \widehat{H}}\right) = \sum_{j=1}^d e^{-\beta E_j},
\end{equation}
where the $E_j$ are the energy eigenvalues of $\widehat{H}$. That settles positivity as well, since $e^{-\beta \widehat{H}}$ has only positive eigenvalues. The mean value $\langle \widehat{H}\rangle$ is calculated as follows
\begin{equation}
	\langle \widehat{H}\rangle = \sum_{j=1}^d E_j \frac{e^{-\beta E_j}}{\mathcal{Z}} = -\frac{1}{\mathcal{Z}}\frac{\partial \mathcal{Z}}{\partial \beta}.
\end{equation}
The minus sign in the last equality comes from the differentiation: every term $e^{-\beta E_j}$ brings down a factor $-E_j$. Written more compactly, then, $\langle \widehat{H}\rangle = -\frac{\partial}{\partial \beta}\log(\mathcal{Z})$.

The entropy of $\rho=\frac{e^{-\beta \widehat{H}}}{\mathcal{Z}}$ is calculated as follows
\begin{equation}
	\eqbox{S(\rho) = k\beta\langle\widehat{H}\rangle + k\log(\mathcal{Z}).}
\end{equation}
Define the temperature in analogy with the microcanonical ensemble:
\begin{equation}
	\frac{1}{T} = \frac{\partial S}{\partial \langle\widehat{H}\rangle} = k\beta \Rightarrow \beta = \frac{1}{kT}.
\end{equation}
Hence
\begin{equation}
	\eqbox{\rho = \frac{e^{-\frac{\widehat{H}}{kT}}}{\mathcal{Z}}.}
\end{equation}
The factors $e^{-E_j/kT}$ are called \emph{Boltzmann factors}.

The Boltzmann factors are a shape, not a mere formula: on the left they occupy the same level scheme at three temperatures, from the spike in the ground state to almost the uniform distribution $1/d$. On the right, the mean energy and the entropy meanwhile climb from $0$ up to $(d-1)\Delta/2$ and $k\log(d)$ -- from the pure state to the maximally mixed one.

% fig_vl22_boltzmann.tex — Was der Faktor e^{-beta E_n} mit einem festen
% Niveauschema macht, wenn man die Temperatur dreht: von der Spitze im
% Grundzustand bis in die Gleichverteilung.
%
% WHY THIS FIGURE.  The lecture derives rho = e^{-beta H}/Z and then
% names e^{-E_j/kT} the Boltzmann factors — but the whole physical
% content of that formula is a *shape*, and a shape is what prose is
% worst at.  Two things the text states but never shows:
%   * the same fixed ladder of levels is occupied completely differently
%     at different T, and the two ends of the temperature scale are the
%     two states this course already knows — the pure ground state
%     (S = 0) and the maximally mixed state (S = k log d);
%   * the boxed entropy formula S = k beta <H> + k log(Z) is a curve.
%     Read at T -> 0 it is the third law that closes this lecture:
%     S -> 0.  Read at T -> infinity it saturates at k log d.
% Panel (a) is the microscopic side (which level carries what weight),
% panel (b) the macroscopic side (what a thermometer and a calorimeter
% would report).  That is the contrast the two panels exist for.
%
% GEOMETRY FROM THE FORMULA.  Nothing is drawn by eye; every bar, curve
% and asymptote is evaluated from the closed forms of this lecture on
% one concrete spectrum: the equally spaced ladder of Vorlesung 12,
% truncated to d = 6 levels and counted from the ground state,
%
%       E_n = n Delta,   n = 0,1,...,5,
%
% (equal spacing, because then Z is a finite geometric series and the
% populations are an exact geometric sequence — the constant ratio
% p_{n+1}/p_n = e^{-beta Delta} is what makes the fanning-out of the
% bars readable).  With x = e^{-Delta/kT} = e^{-beta Delta},
%
%       Z    = sum_{n=0}^{5} x^n = (1 - x^6)/(1 - x),
%       p_n  = x^n / Z = e^{-n Delta/kT} (1-x)/(1-x^6),
%       <H>/Delta = (sum_n n x^n)/Z,
%       S/k  = beta <H> + log(Z)          [the boxed formula of the lecture]
%
% and both limits are exact:
%       T -> 0        p_0 -> 1,   <H> -> 0,        S -> 0,
%       T -> infinity p_n -> 1/d, <H> -> (d-1)Delta/2, S -> k log d.
%
% The three temperatures drawn are kT/Delta = 0.3, 1, 3.  Their exact
% populations (the bar lengths) are
%
%   kT/Delta   x        Z        p_0     p_1     p_2     p_3    p_4    p_5
%     0.3   0.035674  1.036994  0.9643  0.0344  0.0012  4.4e-5  2e-6   6e-8
%     1     0.367879  1.578055  0.6337  0.2331  0.0858  0.0315  0.0116 0.0043
%     3     0.716531  3.050300  0.3278  0.2349  0.1683  0.1206  0.0864 0.0619
%   (1/d = 0.1667 is the T -> infinity value of every one of them.)
%
% and the points marked on the curves of (b) are
%
%   kT/Delta      0.3       1        3      -> limit
%   <H>/Delta   0.0370   0.5671   1.5886      2.5
%   S/k         0.1596   1.0233   1.6448      log 6 = 1.7918
%
% The bars are drawn by pgfplots from these formulas (xcomb at samples
% n = 0..5), not typed in; the numbers above are recorded so the plate
% can be checked without re-running LaTeX.
%
% DESIGN DECISIONS
%   * colour is temperature and nothing else, on the ORDERED ramp the
%     templates keep for ordered parameters (tpclevels):
%     spacecadet -> cgblue -> birrenorange, cold to hot.  It was chosen
%     over the obvious blue-orange-red because those three are NOT
%     monotone in luminance: printed grey, cgblue (0.40) and birrenred
%     (0.34) come out the same, and the legend of (a) would then have
%     two identical rows.  This ramp runs 0.15 -> 0.40 -> 0.54, so cold
%     stays darkest in grey as well.  The same three colours mark the
%     same three temperatures on the curves of (b), which is how one
%     reads off what the hot distribution costs in entropy;
%   * the two curves of (b) are therefore NOT on that ramp: they are
%     birrenred, a colour the plate uses nowhere else, and they share
%     it, because they are two quantities of one and the same state —
%     what tells them apart is the dash pattern and the label written
%     on each.  The temperature dots on them carry a white ring so they
%     read against the curve in colour and in grey alike, and the two
%     dots belonging to one temperature are joined by a short stroke in
%     that temperature's colour — without it the pair at kT = 3 Delta,
%     0.06 apart, fuses into a single blob;
%   * within every level of (a) the bars run cold-top to hot-bottom,
%     always in that order, so the grouping can be read without the
%     legend;
%   * the energy levels are neutral grey rules (levellinegrey): they are
%     the reference scheme, the populations are the data;
%   * grey dotted verticals (classicallim) mark the two limits of the
%     probability axis, p = 1/d and p = 1, and are labelled by the
%     limit that produces them, not by their value.  They are reference
%     lines, hence grey, hence the same style as the asymptotes of (b);
%   * a thin exponential envelope p(E) = e^{-beta E}/Z accompanies each
%     bar group, over exactly the range of the spectrum.  It is the
%     continuous function the levels sample, and it carries the same
%     vertical offset as the bars it belongs to, so it
%     passes exactly through their tips (the offset is legibility only —
%     the true energies are the grey rules);
%   * probability runs horizontally and energy vertically, so the level
%     scheme stands the way every other level diagram in these notes
%     does (Vorlesungen 11, 12, 16).

\begin{center}
\begin{tikzpicture}

% ─────────────────────────────────────────────────────────────
%  The formulas of the lecture, once, as macros.  #1 is kT/Delta.
%  \boltzP is p_n with n as the plot variable x (so it doubles as the
%  continuous envelope p(E) with E = x Delta).
% ─────────────────────────────────────────────────────────────
\def\boltzP#1{ exp(-x/#1)*(1-exp(-1/#1))/(1-exp(-6/#1)) }
% Z, <H>/Delta and S/k with kT/Delta as the plot variable x.
\def\partZ{ ((1-exp(-6/x))/(1-exp(-1/x))) }
\def\meanU{ ((exp(-1/x)+2*exp(-2/x)+3*exp(-3/x)+4*exp(-4/x)+5*exp(-5/x))/\partZ ) }
\def\entrS{ ((1/x)*\meanU + ln(\partZ)) }

% ─────────────────────────────────────────────────────────────
%  (a)  The populations on a fixed ladder, at three temperatures
% ─────────────────────────────────────────────────────────────
\begin{axis}[
    name=pop,
    tpcplot,
    width=7.2cm, height=5.2cm,
    grid=none,
    clip=false,
    xmin=0, xmax=1.13, ymin=-0.55, ymax=6.25,
    xtick={0,0.1666667,0.5,1},
    xticklabels={$0$,$1/d$,$0.5$,$1$},
    ytick={0,1,2,3,4,5},
    xlabel={$p_n$},
    ylabel={$E_n/\Delta$},
    every axis y label/.style={at={(axis description cs:-0.11,0.5)},
                               rotate=90, anchor=south, font=\small},
    title={(a) \tpcifenglish{populations}{Besetzungen}
              $p_n\propto e^{-\beta E_n}$},
    legend style={at={(0.99,0.99)}, anchor=north east, font=\scriptsize},
    legend cell align=left,
]

  % The level scheme E_n = n Delta: reference rules, not data.
  \draw[levellinegrey] (axis cs:0,0) -- (axis cs:1.10,0);
  \draw[levellinegrey] (axis cs:0,1) -- (axis cs:1.10,1);
  \draw[levellinegrey] (axis cs:0,2) -- (axis cs:1.10,2);
  \draw[levellinegrey] (axis cs:0,3) -- (axis cs:1.10,3);
  \draw[levellinegrey] (axis cs:0,4) -- (axis cs:1.10,4);
  \draw[levellinegrey] (axis cs:0,5) -- (axis cs:1.10,5);

  % The two limits of the probability axis.
  \draw[classicallim] (axis cs:0.1666667,-0.45) -- (axis cs:0.1666667,5.65);
  \draw[classicallim] (axis cs:1,-0.45)         -- (axis cs:1,5.65);

  % Envelopes p(E) = e^{-beta E}/Z, each on its own group's offset.
  \addplot[spacecadet,   line width=0.6pt, forget plot,
           domain=0:5, samples=140] ({\boltzP{0.3}}, {x+0.26});
  \addplot[cgblue,       line width=0.6pt, forget plot,
           domain=0:5, samples=140] ({\boltzP{1}},   {x});
  \addplot[birrenorange, line width=0.6pt, forget plot,
           domain=0:5, samples=140] ({\boltzP{3}},   {x-0.26});

  % The populations themselves, sampled at n = 0,...,5.
  \addplot[xcomb, mark=none, spacecadet,   line width=3.4pt,
           domain=0:5, samples=6] ({\boltzP{0.3}}, {x+0.26});
  \addlegendentry{$kT=0.3\,\Delta$}
  \addplot[xcomb, mark=none, cgblue,       line width=3.4pt,
           domain=0:5, samples=6] ({\boltzP{1}},   {x});
  \addlegendentry{$kT=\Delta$}
  \addplot[xcomb, mark=none, birrenorange, line width=3.4pt,
           domain=0:5, samples=6] ({\boltzP{3}},   {x-0.26});
  \addlegendentry{$kT=3\,\Delta$}

  % What the spectrum is.
  \node[lbl, font=\scriptsize, text=birrengray, anchor=west]
    at (axis cs:0.02,5.95) {$E_n=n\Delta$, $d=6$};

  % The limits, written under the value they force.
  \node[lbl, font=\scriptsize, text=birrengray, anchor=north]
    at (axis cs:0.1666667,-1.22) {$T\to\infty$};
  \node[lbl, font=\scriptsize, text=birrengray, anchor=north]
    at (axis cs:1,-1.22) {$T\to0$};
\end{axis}

% ─────────────────────────────────────────────────────────────
%  (b)  The same state seen from outside: <H> and S against T
% ─────────────────────────────────────────────────────────────
\begin{axis}[
    tpcplot,
    at={(pop.south east)}, anchor=south west, xshift=14mm,
    width=5.6cm, height=5.2cm,
    xmin=0, xmax=6.4, ymin=0, ymax=2.85,
    xtick={0,1,2,3,4,5,6}, ytick={0,1,2},
    xlabel={$kT/\Delta$},
    title={(b) \tpcifenglish{energy and entropy}{Energie und Entropie}},
]
  % Saturation values, the T -> infinity ends of both curves.
  \addplot[levellinegrey, domain=0:6.4, samples=2] {2.5};
  \addplot[levellinegrey, domain=0:6.4, samples=2] {1.7917595};
  \node[lbl, font=\scriptsize, text=birrengray, anchor=south west]
    at (axis cs:0.15,2.5) {$(d{-}1)/2$};
  \node[lbl, font=\scriptsize, text=birrengray, anchor=south west]
    at (axis cs:0.15,1.7917595) {$\log d$};

  % At each of the three temperatures of (a), the two curves are read at
  % once; a short stroke in that temperature's colour joins the pair, so
  % it reads as one temperature with two values.  Needed most at
  % kT = 3 Delta, where the readings differ by 0.06 and two bare dots
  % fuse into a single blob — the stroke turns that near-coincidence
  % into what it is: the two curves about to cross.  Kept short on
  % purpose; a full-height vertical would run through the annotations.
  \draw[spacecadet,   line width=0.6pt] (axis cs:0.3,0.00) -- (axis cs:0.3,0.20);
  \draw[cgblue,       line width=0.6pt] (axis cs:1,0.53)   -- (axis cs:1,1.06);
  \draw[birrenorange, line width=0.6pt] (axis cs:3,1.56)   -- (axis cs:3,1.67);

  % <H>/Delta and S/k.  Neutral: these are quantities, not temperatures.
  \addplot[birrenred, line width=1.0pt, domain=0.05:6.4, samples=220]
    {\meanU};
  \addplot[birrenred, line width=1.0pt, densely dashed,
           domain=0.05:6.4, samples=220] {\entrS};

  \node[lbl, anchor=south east] at (axis cs:6.35,2.06)
    {$\langle\widehat{H}\rangle/\Delta$};
  \node[lbl, anchor=north east] at (axis cs:6.35,1.70) {$S/k$};

  % The readings themselves, at the ends of those strokes.
  \node[circle, fill=spacecadet,   inner sep=0.9pt, draw=white, line width=0.4pt] at (axis cs:0.3,0.0370) {};
  \node[circle, fill=spacecadet,   inner sep=0.9pt, draw=white, line width=0.4pt] at (axis cs:0.3,0.1596) {};
  \node[circle, fill=cgblue,       inner sep=0.9pt, draw=white, line width=0.4pt] at (axis cs:1,0.5671) {};
  \node[circle, fill=cgblue,       inner sep=0.9pt, draw=white, line width=0.4pt] at (axis cs:1,1.0233) {};
  \node[circle, fill=birrenorange, inner sep=0.9pt, draw=white, line width=0.4pt] at (axis cs:3,1.5886) {};
  \node[circle, fill=birrenorange, inner sep=0.9pt, draw=white, line width=0.4pt] at (axis cs:3,1.6448) {};

  \node[lbl, font=\scriptsize, text=birrengray, anchor=east, align=right]
    at (axis cs:6.30,0.74)
    {\tpcifenglish{marks: the three}{Marken: die drei}\\
     \tpcifenglish{temperatures of (a)}{Temperaturen aus (a)}};

  % The third law, where this lecture ends.
  \node[lbl, font=\scriptsize, text=birrengray, anchor=west]
    at (axis cs:0.55,0.13) {$T\to0\ \Rightarrow\ S\to0$};
\end{axis}

\end{tikzpicture}
\end{center}


Practically every system in thermal contact with a heat bath is modelled with this \emph{canonical state}, classically as well as quantum mechanically. From now on we also take the following as a definition: \emph{thermal equilibrium} means that the state is given by this equation.

\section*{The Laws of Thermodynamics}

\begin{center}
\parbox[c][5cm][c]{14cm}{\begin{tcolorbox}[colback=banana,arc=0cm,frame hidden,enhanced]{\textbf{Zeroth Law}. If system $A$ is in thermal equilibrium with system $B$, and system $B$ is in thermal equilibrium with system $C$, then system $A$ is in thermal equilibrium with system $C$.}\end{tcolorbox}}
\end{center}
\begin{center}
\parbox[c][5cm][c]{14cm}{\begin{tcolorbox}[colback=banana,arc=0cm,frame hidden,enhanced]{\textbf{First Law}. There is a quantity $U$: the ``internal energy''. The change $\Delta U = Q+W$ of the internal energy equals the heat $Q$ supplied plus the work $W$ done. (Conservation of energy.)}\end{tcolorbox}}
\end{center}
Here \emph{heat} $Q$ is defined as any spontaneous flow of energy from one object to another that stems from a difference in temperature between the two objects. \emph{Work} $W$ is defined in thermodynamics as any other transfer of energy into or out of a system. 
\begin{center}
\parbox[c][5cm][c]{14cm}{\begin{tcolorbox}[colback=banana,arc=0cm,frame hidden,enhanced]{\textbf{Second Law}. There is a quantity ``entropy'', and in adiabatic processes $\Delta S \ge 0$ holds. (Adiabatic means: without exchange of heat with the environment.)}\end{tcolorbox}}
\end{center}
\begin{merksatz}[The perspective of this course]
	Thermodynamics is not fundamental. It is a consequence of quantum mechanics -- a coarse-grained description of the laws of nature for large systems.
\end{merksatz}
\begin{center}
\parbox[c][5cm][c]{14cm}{\begin{tcolorbox}[colback=banana,arc=0cm,frame hidden,enhanced]{\textbf{Third Law}. One cannot fully reach the absolute zero of temperature $T=0$.}\end{tcolorbox}}
\end{center}

This law is not quite right, however, for there are exceptions. It is equivalent to a statement about the entropy: normally the entropy vanishes as $T\to 0$. In some systems -- glass, certain ice crystals, spin glasses -- there is entropy left over at absolute zero; for these the third law does not hold.
 


